\section{Task 2 Code}
\subsection{Main function}
\begin{simplechar}
\begin{lstlisting}
function x = indicatedMethod(Matrix, Vector) % Name of the method as in the textbook
% x stands for obtained result
    [~,Columns] = size(Matrix); % We need to know how big the matrix is in next steps
    % notice the '~', since We assume We use square matrix, We do not need
    % to have another variable for number of rows since it is the same as
    % number of columns
    checkIfMatrixIsSquareMatrix(Matrix);
    [Matrix, Vector] = gaussianEliminationWithPartialPivoting(Columns, Matrix, Vector);
    % Change matrix to upper triangular matrix
    [Matrix, Vector, x] = backSubstitutionPhase(Columns, Matrix, Vector);
    % Get the solution
    x = iterativeResidualCorrection(Matrix, x, Vector); % Improve on the solution
end % end function
\end{lstlisting}
\end{simplechar}

\subsection{checkIfMatrixIsSquareMatrix}
\begin{simplechar}
\begin{lstlisting}
function checkIfMatrixIsSquareMatrix(Matrix)
    [Rows,Columns] = size(Matrix);
    if Rows ~= Columns
        error ('Matrix is not square matrix!');
    end % end if
end % end function
\end{lstlisting}
\end{simplechar}

\newpage
\subsection{gaussianEliminationWithPartialPivoting}
\begin{simplechar}
\begin{lstlisting}
function [Matrix, Vector] = gaussianEliminationWithPartialPivoting(Columns, Matrix, Vector)
    for j = 1 : Columns
            centralElement = max(Matrix(j:Columns,j));
            % We stay in the same row (j) but We change columns, as in the
            % textbook
            [Matrix, Vector] = partialPivoting(Matrix, Vector, j, centralElement, Columns);
            % ensures that a_kk != 0 and reduces errors
            [Matrix, Vector] = gaussianElimination(j, Columns, Matrix, Vector);
            % change matrix into upper triangular matrix
    end % end for
end % end function
\end{lstlisting}
\end{simplechar}

\subsection{partialPivoting}
\begin{simplechar}
\begin{lstlisting}
function [Matrix, Vector] = partialPivoting(Matrix, Vector, j, centralElement, Columns)
    for k = j : Columns
        partialPivotingSwapOneRow(Matrix, Vector, j, k, centralElement);
    end % end for
end % end function
\end{lstlisting}
\end{simplechar}

\newpage
\subsection{partialPivotingSwapOneRow}
\begin{simplechar}
\begin{lstlisting}
function [Matrix, Vector] = partialPivotingSwapOneRow(Matrix, Vector, j, k, centralElement)
    if Matrix(k,j) == centralElement
    swapRowMatrix(Matrix, j, k); % swap jth row with kth row
    swapValueVector(Vector, j, k); % swap jth value with kth value
    end % end if
end % end function
\end{lstlisting}
\end{simplechar}


\subsection{swapRowMatrix}
\begin{simplechar}
\begin{lstlisting}
function Matrix = swapRowMatrix(Matrix, j, k)
    temp =  Matrix(j , :); % ' : ' denote "all elements in jth row"
    Matrix(j , :) = Matrix(k, :);
    Matrix(k, :) = temp; % temp equal to previous value of jth row
end
\end{lstlisting}
\end{simplechar}

\subsection{swapValueVector}
\begin{simplechar}
\begin{lstlisting}
function Vector = swapValueVector(Vector, j, k)
    temp = Vector(j);
    Vector(j) = Vector(k);
    Vector(k) = temp;  % temp equal to previous value of k element of vector
end % end function
\end{lstlisting}
\end{simplechar}

\newpage
\subsection{gaussianElimination}
\begin{simplechar}
\begin{lstlisting}
function [Matrix, Vector] = gaussianElimination(j, Columns, Matrix, Vector)
    for i = j + 1 : Columns
        rowMultiplier = Matrix(i,j) / Matrix(j,j);
        [Matrix, Vector] = substractRows(Matrix, Vector, i, rowMultiplier, j, Columns);
    end % end for
end % end function
\end{lstlisting}
\end{simplechar}

\subsection{substractRows}
\begin{simplechar}
\begin{lstlisting}
function [Matrix, Vector] = substractRows(Matrix, Vector, i, rowMultiplier, j, Columns)
    Vector(i) = Vector(i) - rowMultiplier * Vector(j);
    for curentColumn = 1 : Columns
        Matrix(i,curentColumn) = Matrix(i,curentColumn) - rowMultiplier * Matrix(j, curentColumn);
    end % end for
end % end function
\end{lstlisting}
\end{simplechar}

\newpage
\subsection{backSubstitutionPhase}
\begin{simplechar}
\begin{lstlisting}
function [Matrix, Vector, x] = backSubstitutionPhase(Columns, Matrix, Vector)
    for k = Columns : -1 : 1
    % Start at final column and move by -1 each iteration until we reach 1
        equation = 0;
        for j = k+1 : Columns
            equation = equation + Matrix(k,j) * x(j, 1);
            % even though x is a vector we still need to put '1' to ensure
            % that number of columns in the first matrix matches number of
            % rows in second matrix
        end % end for

        x(k, 1) = (Vector(k,1) - equation) / Matrix(k,k);
        % even though x is a vector we still need to put '1' to ensure
        % that we do not exceed array bounds
    end % end for
end % end function
\end{lstlisting}
\end{simplechar}

\subsection{iterativeResidualCorrection}
\begin{simplechar}
\begin{lstlisting}
function x = iterativeResidualCorrection(Matrix, x, Vector)
    residuum = Matrix*x - Vector; % as in the book
    newResiduum = residuum;
    x = improveSolution(x, newResiduum, residuum, Matrix, Vector);
end % end function
\end{lstlisting}
\end{simplechar}

\newpage
\subsection{improveSolution}
\begin{simplechar}
\begin{lstlisting}
function x = iterativeResidualCorrection(A, x, b)
    r = A*x - b;
    for i = 1 : 100
        [~, ~, deltaX] = solveSystem(A, r);
        newX = x - deltaX;
        r = A*newX - b;
        x = newX;
    end

end % end function

\end{lstlisting}
\end{simplechar}

\newpage
\hypertarget{function2_plot}{\subsection{plotErrorsGaussian}}
\begin{simplechar}
\begin{lstlisting}
function plotErrorsGaussian(maxMatrixSize)

    errorsA = zeros(maxMatrixSize);
    errorsB = zeros(maxMatrixSize);
    errorsAR = zeros(maxMatrixSize);
    errorsBR = zeros(maxMatrixSize);
    for i = 1 : maxMatrixSize
        [~, errorBeforeResidualCorrection, errorAfterResidualCorrection] = indicatedMethod(matrixA(i), vectorA(i));
        errorsA(i) = errorBeforeResidualCorrection;
        errorsAR(i) = errorAfterResidualCorrection;
        [~, errorBeforeResidualCorrection, errorAfterResidualCorrection] = indicatedMethod(matrixB(i), vectorB(i));
        errorsB(i) = errorBeforeResidualCorrection;
        errorsBR(i) = errorAfterResidualCorrection;
    end
    nexttile
    plot(errorsA, '.');
    title('Errors before residual correction for task 2a:');
    xlabel('Size of matrix A');
    ylabel('Errors');
    nexttile
    plot(errorsAR, '.');
    title('Errors after residual correction for task 2a:');
    xlabel('Size of matrix A');
    ylabel('Errors');
    nexttile
    plot(errorsB, '.');
    title('Errors before residual correction for task 2b:');
    xlabel('Size of matrix A');
    ylabel('Errors');
    nexttile
    plot(errorsBR, '.');
    title('Errors after residual correction for task 2b:');
    xlabel('Size of matrix A');
    ylabel('Errors');
end
\end{lstlisting}
\end{simplechar}


